\section{Collisional plasmas}
\label{sec:plasma}

Each of the three facts of Section~\ref{sec:easy} fails in a plasma. This
section covers the particle-to-kinetic closure and the collisional
kinetic-to-fluid closure. The collisionless and magnetized cases follow in
Sections~\ref{sec:collisionless} and~\ref{sec:mag}. A standard reference
for the collisional theory is Helander and Sigmar \cite{HelanderSigmar02}.

\subsection{Parameters}

For electrons of density $n$ and temperature $T$ (energy units),
\[
  \lambda_D=\sqrt{\frac{\varepsilon_0T}{ne^2}},\qquad
  \omega_p=\sqrt{\frac{ne^2}{\varepsilon_0m_e}},\qquad
  \Lambda=\frac{4\pi}{3}\,n\lambda_D^3,\qquad
  \Gamma=\frac{e^2}{4\pi\varepsilon_0\,a\,T},\quad
  a=\Bigl(\frac{3}{4\pi n}\Bigr)^{1/3}.
\]
$\Lambda$ is the number of particles in a Debye sphere and $\Gamma$ is the
ratio of the potential energy of neighbours to the thermal energy. They are
one parameter: $\Gamma=\tfrac13\Lambda^{-2/3}$
(Exercise~\ref{ex:gamma}). A weakly coupled plasma has $\Lambda\gg1$ and
$\Gamma\ll1$.

\subsection{Particles to kinetic: mean field first, collisions second}

The Coulomb force is long-range, so each particle interacts with many
others at once. When $\Lambda\gg1$ the leading-order closure of BBGKY is a
mean field and not a collision term:
\begin{equation}
  \partial_t f_s + \bv\cdot\nabla f_s
  + \frac{q_s}{m_s}\left(\bE + \bv\times\bB\right)\cdot\nabla_{\bv} f_s
  = \sum_{s'}C_{ss'}[f_s,f_{s'}],
  \label{eq:vlasov-landau}
\end{equation}
with $\bE$ and $\bB$ determined self-consistently by Maxwell's equations.
Collisions enter as a correction of relative size $\ln\Lambda/\Lambda$.
They are the cumulative effect of many small-angle deflections, so the
operator is of Fokker--Planck type. For like particles the Landau operator
\cite{Landau36} is
\begin{equation}
  C[f,f]=\frac{q^4\ln\Lambda}{8\pi\varepsilon_0^2m^2}\;
  \nabla_{\bv}\cdot\!\int\mathbf U(\bv-\bv')\cdot
  \bigl[f(\bv')\nabla_{\bv}f(\bv)-f(\bv)\nabla_{\bv'}f(\bv')\bigr]\dd^3v',
  \qquad
  \mathbf U(\bw)=\frac{|\bw|^2\mathbf I-\bw\bw}{|\bw|^3}.
\end{equation}
The logarithm comes from cutting off the Coulomb interaction at the Debye
length and at the distance of closest approach. The Lenard--Balescu
operator \cite{Lenard60,Balescu60} removes the long-range cutoff by
including the dielectric response of the plasma, and is obtained from BBGKY
by neglecting $g_3$ and solving the equation for $g_2$. The expansion
parameter for the particle-to-kinetic closure is weak coupling, not low
density. A rigorous derivation of these equations from particle dynamics is
open; the strategy explored in \texttt{hani-to-algorithms/} is one
approach.

When $\Gamma\gtrsim1$, as for the dust component of a dusty plasma, pair
correlations are not small and this closure fails
(Section~\ref{sec:strong}).

\subsection{Collisions weaken with energy}
\label{sec:coulomb-tail}

The Coulomb collision frequency falls steeply with speed. For electrons
scattering on ions of charge $Z$,
\begin{equation}
  \nu_{ei}(v)=\frac{n_iZ^2e^4\ln\Lambda}{4\pi\varepsilon_0^2m_e^2\,v^3},
  \qquad
  \lambda(v)=\frac{v}{\nu(v)}\propto v^4 .
\end{equation}
There is no single mean free path. The Chapman--Enskog calculation still
goes through (Spitzer and H\"arm \cite{SH53}; Braginskii
\cite{Braginskii65}) and gives $\bq=-\kappa\nabla T$ with
$\kappa\propto T^{5/2}/(Z\ln\Lambda)$ times a function of $Z$. But the
integrand of the heat flux weights high speeds heavily: the heat is carried
by electrons at several times the thermal speed, whose mean free path is
two orders of magnitude longer than the thermal one. The local closure
therefore fails when the thermal mean free path is still of order one
percent of the temperature gradient length or less. In the linear response
language of \eqref{eq:deficit}, the coefficient $K$ is large.

The one-dimensional model results in this repository isolate the mechanism
(Section~\ref{sec:between}). With a velocity-independent collision
frequency the local closure fails only at $k\lambda$ of order one. A
$|v|^{-3}$ tail in the collision frequency moves the breakdown to
$k\lambda\sim10^{-2}$.

\subsection{Repairs used in practice}

\begin{itemize}
  \item \textbf{Flux limiter.} Cap the local flux at a fraction of the
    free-streaming flux,
    \begin{equation}
      q=\min\bigl(q_{\mathrm{SH}},\;f\,nT\vt\bigr)
      \quad\text{or}\quad
      \frac1q=\frac1{q_{\mathrm{SH}}}+\frac1{f\,nT\vt},
    \end{equation}
    with $f$ tuned between about $0.03$ and $0.15$ against experiments or
    kinetic simulations. It is local, it has no predictive value for $f$,
    and it cannot produce preheat ahead of a temperature front.
  \item \textbf{Convolution kernels.} Luciani, Mora and Virmont
    \cite{LMV83} write
    \begin{equation}
      q(x)=\int W(x,x')\,q_{\mathrm{SH}}(x')\,\dd x',
    \end{equation}
    with an exponential kernel whose range is a few tens of thermal mean
    free paths. In a uniform background an exponential kernel of range
    $\ell$ is the multiplier $1/(1+k^2\ell^2)$, which has the form
    \eqref{eq:rational}.
  \item \textbf{Fits to linear kinetic theory.} Epperlein and Short
    \cite{ES91} built a nonlocal model on linear Fokker--Planck results.
    The fit usually quoted from that line of work is
    $\kappa/\kappa_{\mathrm{SH}}=1/\bigl(1+(30\,k\lambda_e)^{4/3}\bigr)$,
    with $\lambda_e=\sqrt{\lambda_{ee}\lambda_{ei}}$ a delocalization length
    and not the thermal mean free path. The formula was confirmed from a
    later paper that quotes it, not from the original.
  \item \textbf{Multigroup diffusion.} Schurtz, Nicola\"i and Busquet
    \cite{SNB00} solve one screened diffusion equation per electron energy
    group and sum the contributions. This is the form used in
    radiation-hydrodynamics codes, because it works in several dimensions
    and avoids an explicit convolution.
  \item \textbf{Moment-based integral closures} for heat flux, viscosity
    and friction together at arbitrary collisionality \cite{JiHeld14}.
\end{itemize}
Comparisons against kinetic solvers show that the ranking of these models
depends on the drive, the temperature contrast and the collision physics
\cite{Brodrick17}. The heat-flux manuscript in this repository proves that
for a thermal drive the exact linear closure is a positive superposition of
screened diffusion responses (a Stieltjes function of $k^2$), which is the
structure the multigroup models assume; it does not tell how many groups
are needed or what their weights are (Section~\ref{sec:repo}).

\subsection{Two fluids: more to close}

A plasma has at least two species with a mass ratio of 1836 or more, which
can have different temperatures and flows, and electromagnetic fields that
are dynamical variables. The fluid equations for each species contain, in
addition to $\boldsymbol\pi_s$ and $\bq_s$, the collisional exchange of
momentum $\mathbf R$ and of heat between species. Braginskii's closure
\cite{Braginskii65} for the unmagnetized direction gives, for $Z=1$,
\begin{equation}
  \mathbf R_\parallel=-0.51\,\frac{m_en}{\tau_e}\,u_\parallel
      -0.71\,n\nabla_\parallel T_e,
  \qquad
  q_{e\parallel}=0.71\,nT_e\,u_\parallel
      -3.16\,\frac{nT_e\tau_e}{m_e}\nabla_\parallel T_e ,
\end{equation}
where $u_\parallel$ is the electron flow relative to the ions and $\tau_e$
the electron collision time. The thermal force (second term of
$\mathbf R$) and the convective heat flux (first term of $q$) carry the
same coefficient $0.71$. That is an Onsager symmetry, a structural
constraint that any closure of the coupled system should respect. The
electron momentum equation with $\mathbf R$ closed is the generalized Ohm's
law (Section~\ref{sec:mag}).

The Debye length, the skin depths and the gyroradii of each species are
distinct scales, and a closure valid at one may fail at another.

\subsection{Closures in sequence}

A plasma model stacks several closures, each with its own defect:
\begin{center}
\renewcommand{\arraystretch}{1.3}
\begin{tabularx}{\textwidth}{@{}P{4.2cm}P{4.2cm}Y@{}}
\toprule
Transition & Small parameter & Fails when \\
\midrule
Particles $\to$ kinetic & $1/\Lambda$ (weak coupling)
   & strong coupling, dense or dusty plasmas \\
Kinetic $\to$ fluid, collisional & $\lambda/L$
   & steep gradients, fast particles \\
Kinetic $\to$ fluid, magnetized & $\rho/L$
   & unmagnetized regions, parallel dynamics \\
Kinetic $\to$ fluid, collisionless & none; matched to linear response
   & nonlinear trapping, strong drives \\
Fluid $\to$ mean transport & none in general
   & always a model (turbulence) \\
\bottomrule
\end{tabularx}
\end{center}
